\chapter{Jesus Cristo desceu aos infernos, ressuscitou dos mortos no terceiro dia}
\chaptermark{Jesus Cristo ressuscitou}

O quinto artigo da Profissão de Fé anuncia o cumprimento glorioso do Mistério Pascal: ``Jesus Cristo desceu aos infernos, ressuscitou dos mortos no terceiro dia''. Este artigo completa a proclamação da vitória de Cristo sobre a morte e o pecado. A descida aos infernos manifesta que Jesus assumiu plenamente a condição humana, chegando até ao último limite da morte. A Ressurreição, por sua vez, é a verdade central da fé cristã, o acontecimento que dá sentido a toda a vida e missão de Jesus. O Catecismo da Igreja Católica, nos parágrafos 632 a 658, apresenta estes mistérios com profundidade e clareza, mostrando como a descida aos infernos e a Ressurreição são inseparáveis na obra da salvação.

\section{Cristo desceu aos infernos (632--637)}

A expressão ``desceu aos infernos'' do Credo não significa que Jesus tenha sofrido as penas dos condenados. Significa, antes, que Ele, após a sua morte na cruz, foi realmente ao reino dos mortos. Jesus experimentou a morte como todos os seres humanos. Na sua alma unida à Pessoa divina, Ele desceu à morada dos defuntos, chamada Sheol ou Hades.

Ao descer aos infernos, Cristo não foi como um prisioneiro, mas como Salvador e Proclamador da Boa Nova. Ele anunciou a salvação aos espíritos que ali se encontravam. Os justos do Antigo Testamento, que haviam morrido na esperança do Redentor, foram libertados por Ele. Cristo abriu as portas do Céu para todos os que tinham aguardado a sua vinda com fé e justiça.

A descida aos infernos manifesta a extensão da missão redentora de Jesus. O alcance da salvação não se limita aos que viveram após a sua vinda, mas estende-se a todos os homens de todos os tempos. ``O Evangelho foi anunciado também aos mortos''. Cristo foi até ao mais profundo da condição humana para vencer a morte na sua própria raiz.

Esta verdade é consoladora para a fé cristã. Jesus não apenas morreu por nós, mas desceu até onde o pecado e a morte haviam lançado a humanidade. Ele foi procurar Adão, o nosso primeiro pai, como a ovelha perdida. Com a sua presença, iluminou as trevas do mundo dos mortos e anunciou a libertação definitiva.

A expressão ``desceu aos infernos'' confirma que Jesus realmente morreu e, pela sua morte, venceu a morte e aquele que tem o poder da morte, isto é, o demónio. Pela sua descida, Ele abriu as portas do Céu aos justos que o haviam precedido. Esta verdade completa a confissão da realidade da morte de Cristo e da universalidade da sua redenção.

\section{No terceiro dia ressuscitou dos mortos (638--658)}

A Ressurreição de Jesus é a verdade central da nossa fé. Os Apóstolos pregaram desde o início que Deus ressuscitou Jesus ao terceiro dia, conforme as Escrituras. Esta verdade coroa toda a vida de Cristo e constitui o fundamento da fé cristã. Sem a Ressurreição, a nossa fé seria vã. Com ela, a esperança da vida eterna torna-se certa.

A Ressurreição é um acontecimento real, historicamente verificável. O túmulo vazio e as aparições do Ressuscitado aos discípulos constituem os dois grandes sinais. As santas mulheres foram as primeiras a descobrir o túmulo vazio e a encontrar o Senhor vivo. Depois, Jesus apareceu a Pedro, aos Doze e a mais de quinhentas pessoas ao mesmo tempo. Estas aparições não foram fruto de alucinações ou de uma fé exaltada, mas de um encontro real com o Senhor vivo.

Os discípulos, inicialmente desanimados e incrédulos, foram transformados pelo encontro com o Ressuscitado. O próprio Jesus lhes censurou a falta de fé e a dureza de coração. A fé pascal nasceu da experiência direta do Cristo vivo, sob a ação da graça divina. A Ressurreição não foi uma invenção dos Apóstolos, mas o acontecimento que os converteu em testemunhas corajosas.

O corpo ressuscitado de Jesus não é um simples regresso à vida terrena, como aconteceu com Lázaro ou com a filha de Jairo. Trata-se de uma vida nova, gloriosa, totalmente penetrada pelo poder do Espírito Santo. O corpo de Cristo ressuscitado possui novas propriedades: não está limitado pelo espaço nem pelo tempo, pode tornar-se presente onde e quando Ele quiser. Ao mesmo tempo, é o mesmo corpo que foi crucificado, pois conserva as marcas da Paixão.

A Ressurreição de Cristo é um acontecimento transcendente. Ninguém foi testemunha direta do próprio momento da Ressurreição. Os evangelistas não descrevem o instante em que Jesus ressuscitou. Trata-se de um mistério que ultrapassa a história, embora tenha deixado sinais históricos bem concretos. O Cristo ressuscitado não se manifesta ao mundo, mas aos seus discípulos, que se tornam testemunhas perante o povo.

A Ressurreição é obra comum da Santíssima Trindade. O Pai manifesta o seu poder ressuscitando o Filho. O Filho retoma a vida que tinha entregue livremente. O Espírito Santo é o poder que vivifica o corpo de Cristo. Esta ação trinitária manifesta a unidade divina na obra da salvação.

A Ressurreição de Cristo é princípio e fonte da nossa própria ressurreição futura. Cristo ressuscitado tornou-se ``as primícias dos que morreram''. Porque Ele vive, nós também viveremos. A sua Ressurreição garante que a morte não tem a última palavra. A vitória de Cristo é a nossa vitória.

O Cristo ressuscitado é o Senhor da história e da criação. Ele detém as chaves da morte e do Hades. Diante do seu nome, todo o joelho se dobra nos céus, na terra e debaixo da terra. A fé na Ressurreição enche o cristão de alegria pascal e de confiança missionária, mesmo no meio das dificuldades da vida presente.

A Ressurreição de Jesus é o acontecimento que dá pleno sentido à sua descida aos infernos. Aquele que desceu até ao mais profundo da morte é o mesmo que ressuscitou glorioso, abrindo o caminho da vida para todos os que creem n'Ele. Esta verdade pascal é o coração da mensagem cristã e a fonte da nossa esperança.
